\documentclass[a4paper, oneside]{article}
\usepackage{graphicx}
\usepackage{amsthm}
\usepackage{amsmath}
\usepackage{amssymb}
\usepackage[a4paper,
            bindingoffset=0.2in,
            left=2cm,
            right=2cm,
            top=2cm,
            bottom=2cm,
            footskip=.25in]{geometry}
\usepackage[italian]{babel}
\usepackage{pgfplots}
\usepackage{tabularx}
\usepackage{tikz}
\usepackage{wrapfig}
\usepackage{color}
\usepackage[d]{esvect}
\definecolor{page}{rgb}{0.129,0.157,0.212}
\pagecolor{page}
\color{white}
\graphicspath{ {./images/} }
\usetikzlibrary{shapes.geometric}
\usetikzlibrary{datavisualization}
\usetikzlibrary{datavisualization.formats.functions}
\usetikzlibrary{patterns}
\pgfplotsset{width=10cm,compat=1.18}

\title{Appunti di Lab II}
\author{Tommaso Miliani}
\date{04-03-26}

\begin{document}
\newtheoremstyle{theoremEnv}
                {}          % Space above
                {}          % Space below
                {\slshape}  % Body font
                {}          % Indent amount
                {\bfseries} % Head font
                {.}         % Punctuation after head
                {\newline}  % Space after theorem head
                {}          % Theorem head spec
\theoremstyle{theoremEnv}

\newtheorem{definition}{Definizione}[section]
\newtheorem{theorem}{Teorema}[section]
\newtheorem{lemma}{Proposizione}[section]
\newtheorem{observation}{Osservazione}[section]
\newtheorem{corollary}{Corollario}[theorem]
\newtheorem{example}{Esempio}[section]
\newtheorem{remark}{Enunciato}[section]

\maketitle

\section{Prima esperienza}
Si è schematizzato quello che accade all'interno del generaetore
con $\vv{E_l}$. All'equilibrio se non si muovono cariche, allora
$\vv{E_l} = - \vv{E_S}$. Si è definita la $fem$ come 
\begin{gather*}
    fem = \int_{B}^{A} \vv{E_l} \cdot \vv{dl}   = \int_{B}^{A} -\vv{E_S} \cdot \vv{dl} = \int_{A}^{B} \vv{E_S} \cdot \vv{dl}     
\end{gather*}   
Si sa fare solo all'equilibrio poiché il campo $\vv{E_S}$ è equivalente a quello 
non conservativo $\vv{E_l}$. Dato che è conservativo $\vv{E_S}$, allora integrale si può
eguagliare ad un potenziale. DUnque si definisce operativamente
\begin{gather*}
    fem = f = \epsilon = \left.(V_A - V_B)_m \right|_{i = 0 }
\end{gather*}
Sperimentalemente si può misurare la differenza di potenziale e si trova che, 
in funzione della corrente, il grafico non è costante ma alla corrente 
massima la tensione crolla. QUesto si spiega perché la corrente che passa fuori dal
circuito deve anche passare all'interno del generatore, essendo esso stesso 
un conduttore, avrà una certa resistenza. Per spiegare dunque questo andamento, 
si può rappresentare la caduta di potenziale come il valore di resistenza finito (anche
se piccola) del generatore stesso. Ecco come, in tutti gli altri casi, 
a circuito chiuso, la corrente $i$ che passa non è infinita e la differenza di potenziale
ai cavi sarà uguale alla 
\begin{gather*}
    (V_A - V_B)_m = fem - ir
\end{gather*}
Questo vale fino a che non ci si avvicina alla corrente  $i_{max}$ erogabile 
dal generatore. 
\paragraph{Metodo Potenziometrico} \mbox{} \\
Si possono realizzare circuiti con due maglie prendendo un generatore stabile con 
un partitore di tenisoni a cui attacco un amperometro con resistenza e e generatore extra incognito.
Il circuito della prima maglia permette di variare in maniera continua la tensione 
ad uno dei capi. Se nel divisore di tensione si riesce a tarare il rapporto giusto di resistenze, allora
attraverso la maglia che contiene l'amperometro non passerà più corrente. Questo 
prende il nome di \textbf{metodo potenziometrico}. Questa misura non si può effettuare 
da sola, infatti la misura in laboratorio sarà composta da due step:
\begin{enumerate}
    \item  La prima parte della misura è di taratura: indipendentemente dalla
    fem del primo generatore, conosciuto il rapporto di tensione nel divisore, si utilizza
    un generatore di tensione ideale $\epsilon_R$ e $r_R$, da regolare in modo tale che l'amperometro segni zero.
    Questa taratura dipenderà dal rapporto nel divisore di tensione. 
    \item Si collega poi ad un secondo circuito con resistenza
    e fem ignota e, tramite la misura tra il rapporto $r_R$ e $r_x$ nei 
    due divisori di tensioni si riesce a trovare il rapporto tra $\epsilon_R$ e $\epsilon_x$
\end{enumerate}
Questi sistemi riescono a determinare la tensione e la resistenza con una precisione 
fino a
\begin{gather*}
    \frac{\Delta V}{V} \sim 10^{-9} \qquad \frac{\Delta R}{R} \sim 10^{-9}
\end{gather*}

\clearpage
\section{Risoluzione di esercizi}
I campi non conservativi non compaiono all'interno della legge delle 
maglie di Kirkhoff. A partire dall'equazione locale (che vale
punto per punto per un circuito):
\begin{gather*}
    \vv{E} = \rho_R \cdot  \vv{J}  
\end{gather*}
In realtà in ogni punto del circuito, anche se vale punto per punto, non è 
detto che tutti i punti siano alla stessa tensione e alla stessa resistenza.
Si può però pensare che in un certo punto del circuito ci sia della
resistenza concentrata. Come si può esprimere l'integrale tra $A$ e $B$:
\begin{gather*}
    \int_{A}^{B} \vv{E} \cdot \vv{dl} = \int_{A}^{B} \rho_R \vv{J} \cdot \vv{dl}      
\end{gather*}
questo campo $\vv{E}$ è la somma dei campi del generatore più il campo elettrostatico
conservativo che compare a causa dell'accumulo di carica:
\begin{gather*}
    \vv{E} = \vv{E_e} + \vv{E_s} \ \Longrightarrow \ \int_{A}^{B} \vv{E_S} \cdot \vv{dl} + \int_{A}^{B} \vv{E_e} \cdot \vv{dl} = \int_{A}^{B} \rho_R \vv{J} \cdot \vv{dl}             
\end{gather*} 
CI sarà sicuramente una parte conservativa di campo ed una non conservativa. Tutti i campi
conservativi sono riassunti in $\vv{E_S}$. Facendo la somma lungo il percorso specificato, 
l'integrale sarà dunque la differenza di potenziale tra i due punti (per definizione) 
\begin{gather*}
    V_A - V_B + \int_{-}^{+} \vv{E_{e1}} \cdot \vv{dl} + \int_{+}^{-} \vv{E_e} \cdot \vv{dl} = \int_{A}^{B} \rho_R \vv{J} \cdot \vv{dl}  = R_1 i_1 + R_2 i_2 + R_3 i_3 
\end{gather*} 
Il primo integrale da $f_1$ ed il secondo $-f_2$ il segno di $f_2$ è dovuto al segno di percorrenza
scelta di $\vv{J}$ per utilizzare la legge delle maglie. Dunque la legge di Ohm generalizzata in un ramo è 
\begin{gather*}
    V_A - V_B + \sum_{i} f_i = \sum_{k } iR_k
\end{gather*} 

\begin{example}[Risoluzione di un circuito con due metodi differenti]
    Se si volesse trovare la corrente $i_3$ che scorre lungo il ramo che 
    contiene la resistenza $R_3$:

    Si può applicare le leggi di Kirckhoff ai nodi e alle due maglie 
    per risolvere il circuito e trovare la corrente incognita. Un'altra 
    metodologia è utilizzare il principio di sovrapposizione:
    in un circuito con generatori che spingono la corrente nel circuito, si possono considerare
    la sovrapposizione degli effetti dei due generatori:
    \begin{gather*}
        i_3 = C_1 \epsilon_1 + C_2 \epsilon_2 = i_3' + i_3''
    \end{gather*}
    Il primo contributo è quando $\epsilon_2$ è cortocircuitato e 
    il secondo è quando $\epsilon_1$ viene cortocircuitato. Cortocircuitando 
    il secondo generatore, esso verrà tolto dal suo ramo corrispondente, ora si 
    risolve trovando con la legge dei nodi al nodo $N$ si ha che
    \begin{gather*}
        i_1 = i_2 + i_3'
    \end{gather*}
    Calcolando la resistenza totale (come due resistenze in parallelo in serie con una terza)
    adesso, si può trovare la corrente $i_1$ come
    \begin{gather*}
        R_T = R_1 + \frac{R_2R_3}{R_2 + R_3} \ \Longrightarrow \ i_1 = \frac{\epsilon_1}{R_T}
    \end{gather*}
    Di nuovo, secondo le leggi di Kirckhoff, la corrente si divide in maniera inversamente 
    proporzionale nei circuiti:
    \begin{gather*}
        i_2 R_2 = i_3' R_3 \ \Longrightarrow \ \frac{i_3'}{i_2} = \frac{R_2}{R_3} 
    \end{gather*}
    Dato che
    \begin{gather*}
        i_2 = \frac{R_3}{R_2 + R_3} i_1 \qquad i_3' = \frac{R_2}{R_2 + R_3} i_1 
    \end{gather*}
    Allora 
    \begin{gather*}
        i_3' = \frac{R_2}{R_1R_2 + R_1 R_3 + R_2R_3}\epsilon_1
    \end{gather*}
    Si può applicare lo stesso procedimento per trovare $i_3''$:
    \begin{gather*}
        i_3'' = -\frac{R_1}{R_1R_2 + R_1 R_3 + R_2 R_3} \epsilon_2
    \end{gather*}
    Ottenendo allora 
    \begin{gather*}
        i_3 = \frac{R_2 \epsilon_1 - R_1 \epsilon_2}{R_1R_2 + R_2 R_3+ R_1 R_3}
    \end{gather*}
\end{example}



\end{document}